\unitopen{24}
\begin{candos}
\item napisać prosty życiorys po niemiecku i opowiedzieć go ustnie;
\item opisać swoje doświadczenie zawodowe w Perfekcie i Präsens;
\item przejść test końcowy A2: celownik, Perfekt, zdania podrzędne, stopniowanie, zwroty uprzejme;
\item rozumieć format egzaminu A2 (cztery moduły).
\end{candos}

\section{Słownictwo}
\vocabtable{24}

\section{Życiorys po niemiecku}
Niemiecki \de{Lebenslauf} jest tabelaryczny, odwrotnie chronologiczny (najnowsze na górze) i zwykle ma 1--2 strony. Typowe sekcje:
\de{Persönliche Daten}, \de{Berufserfahrung}, \de{Ausbildung}, \de{Kenntnisse}, \de{Sprachen}, \de{Weiterbildung}.

\begin{center}\small
\begin{tabular}{@{}p{0.25\linewidth}p{0.65\linewidth}@{}}
\toprule
\de{seit 03/2022} & \de{Softwareentwickler, Firma Alpha, Warschau}\\
 & \de{Entwicklung von Web-APIs mit .NET, Betreuung von drei Junior-Entwicklern}\\
\de{2018--2022} & \de{Entwickler, Firma Beta, Krakau}\\
\de{2012--2017} & \de{Studium der Informatik, Technische Universität Krakau}\\
 & \de{Abschluss: Master of Science}\\
\bottomrule
\end{tabular}
\end{center}

\begin{zasada}[Czasy w CV i w rozmowie]
W tabeli CV zdania są często rzeczownikowe (bez czasownika). W rozmowie o doświadczeniu: Perfekt dla zakończonych działań
(\de{Ich habe drei Jahre bei Beta gearbeitet}), Präsens dla trwającego (\de{Ich arbeite seit 2022 bei Alpha}).
\end{zasada}

\begin{falle}
Polskie CV zawiera często „zainteresowania” i fotografię; w Szwajcarii zdjęcie bywa wciąż spotykane, ale nie
wymagane \verify{(praktyka zależy od branży i firmy)}. Błąd językowy: \de{Ich habe studiert Informatik} (szyk). Poprawnie:
\de{Ich habe Informatik \textbf{studiert}}. I tłumaczenie „studiowałem” jako \de{Ich studierte} w CV ustnym, gdy
w mowie powinno być Perfekt.
\end{falle}

\begin{szwajcaria}
W Szwajcarii pracodawcy cenią \de{Arbeitszeugnisse} (świadectwa pracy) i \de{Referenzen} (kontakty do polecających).
Poproś poprzedniego pracodawcę o świadectwo w języku, w którym możesz się z nim posługiwać. Przy rekrutacji z zagranicy firmy
akceptują angielskie CV, ale znajomość niemieckiego jest atutem, o ile zgadza się z deklaracją.
\end{szwajcaria}

\begin{praca}
\de{In meiner letzten Position habe ich ein Team von drei Entwicklern geleitet.} \de{Ich habe eine REST-API mit .NET
entwickelt.} \de{Ich bin verantwortlich für die Cloud-Infrastruktur.} \de{Ich spreche fließend Englisch und lerne seit einem Jahr Deutsch.}
\end{praca}

\section{Dialogi}
\textbf{Dialog 1 --- pierwsza rozmowa}
\begin{dialog}
\turn{A}{Erzählen Sie bitte etwas über sich.}{Proszę opowiedzieć coś o sobie.}
\turn{B}{Ich heiße Piotr Nowak und komme aus Polen. Ich bin Softwareentwickler und arbeite seit vier Jahren mit .NET.}{Nazywam się Piotr Nowak i pochodzę z Polski. Jestem programistą i pracuję z .NET od czterech lat.}
\turn{A}{Was haben Sie studiert?}{Co Pan studiował?}
\turn{B}{Ich habe Informatik in Krakau studiert und 2017 meinen Abschluss gemacht.}{Studiowałem informatykę w Krakowie i w 2017 ukończyłem studia.}
\turn{A}{Und warum möchten Sie in der Schweiz arbeiten?}{A dlaczego chce Pan pracować w Szwajcarii?}
\turn{B}{Weil ich neue Erfahrungen sammeln möchte und die Firmen hier sehr professionell sind.}{Ponieważ chcę zdobyć nowe doświadczenia, a firmy tutaj są bardzo profesjonalne.}
\end{dialog}

\section{Test końcowy A2}
Wynik poniżej 80\% (10 z 12) oznacza powtórkę tygodni 14, 16 i 18.

\begin{cwiczenia}
\exitem{Celownik: \de{Ich helfe \dots{} (er)}.}{\de{ihm}}
\exitem{Perfekt: \de{Ich / nach Zürich / fahren}.}{\de{Ich bin nach Zürich gefahren.}}
\exitem{Perfekt: \de{Wir / das Release / planen}.}{\de{Wir haben das Release geplant.}}
\exitem{Zdanie podrzędne: \de{Ich bleibe zu Hause, weil ich \dots{} (krank sein)}.}{\de{krank bin}}
\exitem{Zdanie z \de{wenn}, zaczynając od niego: \de{wenn du Zeit hast / ich rufe dich an}.}{\de{Wenn du Zeit hast, rufe ich dich an.}}
\exitem{Stopniowanie: \de{gut} $\rightarrow$ \dots{} $\rightarrow$ \dots}{\de{besser}, \de{am besten}}
\exitem{Przymiotnik: \de{ein neu\dots{} Laptop}.}{\de{ein neuer Laptop}}
\exitem{Zwrotny: \de{Ich interessiere \dots{} für Musik.}}{\de{mich}}
\exitem{Uprzejma prośba o pomoc.}{\de{Könnten Sie mir bitte helfen?}}
\exitem{Początek formalnego e-maila do pani Keller.}{\de{Sehr geehrte Frau Keller,}}
\exitem{Präteritum: \de{Ich \dots{} gestern krank.}}{\de{war}}
\exitem{Przetłumacz: „Pracuję jako programista od czterech lat.”}{\de{Ich arbeite seit vier Jahren als Softwareentwickler.}}
\end{cwiczenia}

\section{Trening słownictwa (generowany)}
\drill{24}

\weekend{\textbf{Zadanie końcowe A2.} Wygłoś trzyminutową prezentację o sobie: życiorys, doświadczenie, plany i dlaczego Szwajcaria.
Nagraj, wyślij do korekty i porównaj z nagraniem A1. Format egzaminu A2 (np.\ Goethe-Zertifikat A2) to zwykle
czytanie, słuchanie, pisanie i mówienie; przejrzyj arkusz przykładowy ze strony egzaminatora.}
